% Self-contained section fragment for integration; ordinary amsmath/amsthm suffice.
% Define \Li as \operatorname{Li} if the host does not already do so.
\section[A bilateral proof of bicolored parity and the Gaussian Family--2 identity]{A bilateral proof of bicolored parity\\and the Gaussian Family--2 identity}
\label{bip:sec:main}

The canonical manuscript's equation \texttt{double:eq:fam2-example}
was supported by integer-relation discovery and numerical evaluation.
We give a complete proof, and simplify its right-hand side to three
single-value products.  The proof is a concrete application of known
parity theory: Nakamura's unit-circle functional relation
\cite[Proposition~1.2]{bip:Nakamura2012} precedes Panzer's general
multiple-polylogarithm parity theorem \cite[Theorem~1.3 and (3.2)]{bip:Panzer2017}.
Our contribution here is the explicit formula in the manuscript's
\emph{decreasing-index} convention, a direct bilateral proof including
indices equal to one and the collision $xy=1$, and its exact application
to the recorded Gaussian relation.  It is not a claim of priority for
unit-circle parity or for reducibility of all the constants displayed below.

\subsection{Conventions, ordinary convergence, and the finite formula}
Write
\[
 Z_{p,q}(x,y)=\Li_{p,q}(x,y)
 =\sum_{n>m\ge1}\frac{x^ny^m}{n^pm^q},\qquad p,q\ge1.
\]
In this section $|x|=|y|=1$ and $x,y\ne1$; the product $r=xy$ is
allowed to equal one.  These hypotheses include the Gaussian pairs
$(i,i)$, $(-1,i)$, $(i,-1)$, and $(i,-i)$.
The ordinary nested series converges, including when $p=1$.
Indeed, for $y\ne1$, the inner partial sums
$H_{n-1}^{(q)}(y)=\sum_{m<n}y^m/m^q$ are bounded; the sequence
$H_{n-1}^{(q)}(y)/n^p$ tends to zero and its successive differences
are absolutely summable.  Summation by parts against the bounded
partial sums of $x^n$ proves convergence.

For $|u|=1$ put
\[
 \mathcal B_j(u)=\Li_j(u)+(-1)^j\Li_j(u^{-1})\quad(u\ne1).
\]
At $u=1$ define
\[
 \mathcal B_j(1)=\begin{cases}2\zeta(j),&j\text{ even},\\0,&j\text{ odd}.
 \end{cases}
\]
In particular $\mathcal B_1(1)=0$ is a symmetric bilateral value;
it is not an assignment of a finite value to $\Li_1(1)$.
For $u=e^{2\pi i t}$, $0<t<1$, the classical inversion identity gives
\[
 \mathcal B_j(u)=-\frac{(2\pi i)^j}{j!}B_j(t),
\]
where $B_j(t)$ is the Bernoulli polynomial. Thus all these factors are
explicit rational multiples of $(i\pi)^j$ at roots of unity.

For $w=p+q$ define
\begin{align}
 \mathcal R_{p,q}(x,y)={}&(-1)^q\Biggl[
 \sum_{j=1}^q(-1)^j\binom{w-j-1}{p-1}
       \mathcal B_j(r)\Li_{w-j}(x)\notag\\
 &\qquad+\sum_{j=1}^p\binom{w-j-1}{q-1}
       \mathcal B_j(r)\Li_{w-j}(y^{-1})\notag\\
 &\qquad-\binom{w-1}{p}\Li_w(y^{-1})
       -\binom{w-1}{q}\Li_w(x)
       \notag\\
 &\qquad-\mathcal B_p(x)\Li_q(y^{-1})\Biggr]
       +(-1)^w\Li_w(r^{-1}).
 \label{bip:eq:R}
\end{align}
No $\Li_1(1)$ occurs in this expression: the arguments of the possible
order-one single values are $x$ and $y^{-1}$, both different from one;
the last single value has order $w\ge2$.

\begin{theorem}[Explicit bicolored parity, in decreasing-index convention]
\label{bip:thm:parity}
For the parameters above,
\begin{equation}
 \boxed{\quad Z_{p,q}(x,y)-(-1)^wZ_{p,q}(x^{-1},y^{-1})
      =\mathcal R_{p,q}(x,y).\quad}
 \label{bip:eq:parity}
\end{equation}
Consequently
\[
 \Re Z_{p,q}(x,y)=\frac12\mathcal R_{p,q}(x,y)
       \quad(w\text{ odd}),
 \qquad
 \Im Z_{p,q}(x,y)=\frac1{2i}\mathcal R_{p,q}(x,y)
       \quad(w\text{ even}).
\]
At roots of unity these are finite formulas in single polylogarithms
and powers of $\pi$.  The opposite parity component is not evaluated
by this theorem.
\end{theorem}

\subsection{Proof with the summation order made explicit}
For $k\ge1$ form the absolutely convergent bilateral sum
\[
 C_{p,q}(k;r)=\sum_{m\in\mathbb Z\setminus\{0,-k\}}
       \frac{r^m}{(m+k)^pm^q}.
\]
The denominator is $O(|m|^w)$ as $|m|\to\infty$, and $w\ge2$.
The elementary partial-fraction identity
\begin{align}
 \frac1{(m+k)^pm^q}
 ={}&\sum_{j=1}^q(-1)^{q-j}\binom{w-j-1}{p-1}
       \frac1{k^{w-j}m^j}\notag\\
 &+(-1)^q\sum_{j=1}^p\binom{w-j-1}{q-1}
       \frac1{k^{w-j}(m+k)^j}
 \label{bip:eq:pf}
\end{align}
holds away from its poles.  It follows either by matching the principal
parts at $m=0,-k$, after which the difference is a rational function
with no poles and vanishing at infinity, or by induction on $p+q$.

Sum \eqref{bip:eq:pf} over a symmetric cutoff in $m$.  The possible
$j=1$ terms are legitimate in this order:
\[
 \underset{M\to\infty}{\lim}\sum_{0<|m|\le M}\frac{r^m}{m^j}
      =\mathcal B_j(r).
\]
For $r\ne1$, the two one-sided $j=1$ series converge by Dirichlet's
test; for $r=1$ they cancel at every symmetric cutoff.  Translating
the cutoff by the fixed integer $k$ changes the sum by finitely many
terms tending to zero, so the substitution $m+k=\ell$ is valid.
Removing the omitted terms $m=-k$ and $\ell=k$ gives
\begin{align}
 C_{p,q}(k;r)={}&
 \sum_{j=1}^q(-1)^{q-j}\binom{w-j-1}{p-1}
   \frac{\mathcal B_j(r)}{k^{w-j}}\notag\\
 &+(-1)^qr^{-k}\sum_{j=1}^p\binom{w-j-1}{q-1}
   \frac{\mathcal B_j(r)}{k^{w-j}}\notag\\
 &-\frac{(-1)^q}{k^w}
       \left[\binom{w-1}{p}r^{-k}+\binom{w-1}{q}\right].
 \label{bip:eq:C}
\end{align}
The last line uses the elementary hockey-stick sums
$\sum_{j=1}^q\binom{w-j-1}{p-1}=\binom{w-1}{p}$ and its
$p,q$ counterpart.

For $0<\rho<1$ multiply \eqref{bip:eq:C} by $\rho^kx^k$ and sum over
$k\ge1$. All the original double sums are now absolutely convergent.
On the right, let $\rho\uparrow1$ in the resulting single
polylogarithms; their arguments are $x$ and $y^{-1}$, neither equal
to one.  The result is
\begin{align}
 \mathcal A={}&(-1)^q\Biggl[
 \sum_{j=1}^q(-1)^j\binom{w-j-1}{p-1}
       \mathcal B_j(r)\Li_{w-j}(x)\notag\\
 &\quad+\sum_{j=1}^p\binom{w-j-1}{q-1}
       \mathcal B_j(r)\Li_{w-j}(y^{-1})
       -\binom{w-1}{p}\Li_w(y^{-1})
       -\binom{w-1}{q}\Li_w(x)\Biggr].
 \label{bip:eq:A}
\end{align}

To identify the left side in the ordinary nested convention, split
its lattice points, with $n=m+k$, into three disjoint regions:
$n>m>0$, $0>n>m$, and $n>0>m$. They yield, as $\rho\uparrow1$,
\begin{equation}
 \mathcal A=Z_{p,q}(x,y)+(-1)^wZ_{q,p}(y^{-1},x^{-1})
            +(-1)^q\Li_p(x)\Li_q(y^{-1}).
 \label{bip:eq:regions}
\end{equation}
For completeness, the gap Abel factors $\rho^{n-m}$ used here need
not be identified informally with outer-index radial factors.  The
first region is $Z_{p,q}(\rho x,y/\rho)$ and has the integral
\begin{equation}
 Z_{p,q}(a,b)=\frac1{(p-1)!}
    \int_0^1(-\log t)^{p-1}
            \frac{a\Li_q(abt)}{1-at}\,dt.
 \label{bip:eq:integral}
\end{equation}
For $a=\rho x$, $ab=r$, its limit follows by dominated convergence:
$1-\rho xt$ is bounded away from zero uniformly for $\rho$ near one;
$\Li_q(rt)$ is bounded when $q>1$ and at worst logarithmic near
$t=1$ when $q=1,r=1$.  These are integrable majorants.  The same
integral identifies this limit with the ordinary nested series via
outer-index Abel summation.  The negative region is handled
identically, and the cross region is the product
$(-1)^q\Li_p(\rho x)\Li_q(\rho/y)$.  This proves
\eqref{bip:eq:regions} for all stated boundary cases.

Finally the convergent stuffle identity gives
\[
 Z_{q,p}(y^{-1},x^{-1})=
 \Li_q(y^{-1})\Li_p(x^{-1})-\Li_w(r^{-1})
 -Z_{p,q}(x^{-1},y^{-1}).
\]
It follows by multiplying finite partial sums and passing to their
limits; both nested series have already been shown to converge.
Substituting this into \eqref{bip:eq:regions} and using
$\mathcal B_p(x)=\Li_p(x)+(-1)^p\Li_p(x^{-1})$ gives
\eqref{bip:eq:parity}.  Complex conjugation proves the two displayed
real/imaginary consequences.

The collision $r=1$ has therefore been proved directly. One can also
see why it is harmless in \eqref{bip:eq:R}: its $\mathcal B_1(r)$
coefficient is a constant multiple of
$\Li_{w-1}(y^{-1})-\Li_{w-1}(x)$, which vanishes at $r=1$.
The jump in the individual bilateral order-one factor cannot create
a jump in the complete expression.

\subsection{Proof and simplification of the recorded Family--2 identity}
Let $L=\log2$, $G=\beta(2)$, and $v=(1+i)/2$. The exact depth-one
values required at weight three are
\[
 \Li_1(i)=-\frac L2+\frac{i\pi}4,\quad
 \Li_2(i)=-\frac{\pi^2}{48}+iG,\quad
 \Li_3(i)=-\frac{3\zeta(3)}{32}+\frac{i\pi^3}{32},
\]
with $\Li_n(-i)=\overline{\Li_n(i)}$ and
$\Li_n(-1)=-(1-2^{1-n})\zeta(n)$ for $n>1$.
Here $\mathcal B_1(-1)=0$ and
$\mathcal B_2(-1)=-\pi^2/6$. Inserting these numbers in
\eqref{bip:eq:R}, first at $(p,q)=(1,2)$ and then at $(2,1)$,
gives
\begin{align}
 \Re\Li_{1,2}(i,i)
   &=\frac{33}{64}\zeta(3)-\frac\pi4G+\frac{\pi^2L}{24},
       \label{bip:eq:12}\\
 \Re\Li_{2,1}(i,i)
   &=\frac{15}{64}\zeta(3)-\frac{\pi^2L}{32}.
       \label{bip:eq:21}
\end{align}
Consequently the desired difference has the short exact form
\begin{equation}
 \boxed{\quad
 \Re\bigl(\Li_{1,2}(i,i)-\Li_{2,1}(i,i)\bigr)
 =\frac9{32}\zeta(3)-\frac\pi4G+\frac7{96}\pi^2\log2.
 \quad}
 \label{bip:eq:family2-short}
\end{equation}

To identify this with the exact equation printed in the source, write
$\mu_2=\Re\Li_2(v)$, $\lambda_2=\Im\Li_2(v)$ and
$h_{a,b}=\Re\Li_{a,b}(i,1)$. The classical Landen formula
\[
 \Li_2\!\left(\frac z{z-1}\right)
      =-\Li_2(z)-\frac12\log^2(1-z)
\]
(applied along the segment $z=it$, $0\le t\le1$) gives
\[
 \mu_2=\frac{5\pi^2}{96}-\frac{L^2}{8},\qquad
 \lambda_2=G-\frac{\pi L}{8}.
\]
Its derivative is an identity of elementary logarithms and its
constant is fixed at $z=0$, so no numerical recognition is needed.
The shuffle identity immediately gives
\[
 h_{1,1}=\frac{L^2}{8}-\frac{\pi^2}{32}.
\]
For the remaining height-one value put
\[
 S(z)=\Li_{2,1}(z,1)=\frac12\int_0^z\frac{\log^2(1-t)}t\,dt.
\]
The needed inversion is already recorded after the proof of the source's
Proposition~1.12, ``Nielsen inversion with explicit endpoint data''
\cite{SourceManuscript}; both the inversion and its Gaussian specialization
are also explicit in Panzer's discussion following Corollary~1.4
\cite{bip:Panzer2017}.  To make the present proof independent of an
unlabelled source display, let $u=\log(-z)$ on
$\mathbb C\setminus[0,\infty)$.  The identity
$\log(1-1/z)=\log(1-z)-u$ shows by differentiation that
\[
 S(z)+S(1/z)-\Li_3(z)+u\Li_2(z)+\frac{u^3}{6}
\]
is constant.  Here
\[
 S(1)=\frac12\int_0^1\frac{\log^2t}{1-t}\,dt=\zeta(3),
\]
where the last equality follows by integrating the geometric series.
Letting $z\to1$ along the upper unit semicircle gives $u\to-i\pi$;
the two cubic imaginary terms cancel because $\zeta(2)=\pi^2/6$.
The constant is therefore $\zeta(3)$, proving
\begin{equation}\label{bip:eq:nielsen-inversion}
 S(z)+S(1/z)=\Li_3(z)-\log(-z)\Li_2(z)
               -\frac16\log^3(-z)+\zeta(3).
\end{equation}
At $z=i$, conjugation gives $S(1/i)=\overline{S(i)}$ and
$\log(-i)=-i\pi/2$, so
\[
 h_{2,1}=\frac{29}{64}\zeta(3)-\frac\pi4G.
\]
Finally $\Li_1(i)\Li_2(i)=\Li_{1,2}(i,1)+2\Li_{2,1}(i,1)$
gives
\[
 h_{1,2}=\Re\bigl(\Li_1(i)\Li_2(i)\bigr)-2h_{2,1}
       =\frac{\pi^2L}{96}+\frac\pi4G-\frac{29}{32}\zeta(3).
\]
Direct rational substitution now yields
\[
 3\zeta(3)-4L\mu_2-\pi\lambda_2-4Lh_{1,1}+3h_{1,2}
 =\frac9{32}\zeta(3)-\frac\pi4G+\frac7{96}\pi^2L.
\]
This proves \texttt{double:eq:fam2-example} precisely as written,
and proves its stronger simplification \eqref{bip:eq:family2-short}.
The statement needs no conjectured independence of the three constants.

\subsection{Mixed-color examples beyond the displayed source equation}
The same finite formula gives the following exact evaluations:
\begin{align}
 \Re\bigl(\Li_{1,2}(-1,i)-\Li_{2,1}(i,-1)\bigr)
   &=\frac{15}{16}\zeta(3)-\frac\pi2G+\frac{\pi^2L}{48},
       \label{bip:eq:mix3}\\
 \Im\bigl(\Li_{1,3}(-1,i)-\Li_{3,1}(i,-1)\bigr)
   &=\frac{3\pi^3L}{32}-3\beta(4)+\frac{21\pi\zeta(3)}{64},
       \label{bip:eq:mix4}\\
 \Re\bigl(\Li_{1,4}(-1,i)-\Li_{4,1}(i,-1)\bigr)
   &=\frac{7\pi^4L}{11520}-\frac\pi2\beta(4)
          +\frac{\pi^2\zeta(3)}{32}+\frac{135}{128}\zeta(5),
       \label{bip:eq:mix5}\\
 \Im\bigl(\Li_{1,5}(-1,i)-\Li_{5,1}(i,-1)\bigr)
   &=\frac{5\pi^5L}{512}-5\beta(6)
          +\frac{465\pi\zeta(5)}{1024}
          +\frac{3\pi^3\zeta(3)}{64}.
       \label{bip:eq:mix6}
\end{align}
Both indices \emph{and} colors are swapped in these antisymmetric
expressions. They are elementary rational specializations of
Theorem~\ref{bip:thm:parity}, not additional parity conjectures.
The accompanying symbolic generator derives the Gaussian formulas at
all requested weights; a separate Mellin-integral evaluator checks the
left sides without reusing the theorem as its evaluation algorithm.

As a collision check, the case $p=q=1$, $(x,y)=(i,-i)$ gives
\[
 \Im\Li_{1,1}(i,-i)=G-\frac{\pi\log2}{8}.
\]
Here the symmetric bilateral value $\mathcal B_1(1)=0$ is essential to
avoid introducing a spurious divergent $\Li_1(1)$.


\subsection{An all-weight family of ordinary logarithmic integrals}
The single-inner-index specialization has an additional cancellation:
no even beta values remain in its real odd-weight component.

\begin{corollary}[Odd logarithmic moments]
For $m\ge1$, put $w=2m+1$ and
\[
 I_m=\frac1{(2m-1)!}\int_0^1
       \frac{x(-\log x)^{2m-1}\log(1+x)}{1+x^2}\,dx.
\]
Then
\begin{align}
 I_m={}&\frac12(1-2^{-2m})
        \left(1-\frac{2m+1}{2^{2m+1}}\right)\zeta(2m+1)\notag\\
 &-\sum_{k=1}^{m-1}
  \frac{(1-2^{1-2k})(1-2^{1-w+2k})}{2^{w-2k}}
       \zeta(2k)\zeta(w-2k)\notag\\
 &-\frac12(1-2^{1-2m})(1-2^{-2m})\zeta(2m)\log2.
 \label{bip:eq:moment-family}
\end{align}
In particular,
\begin{align*}
 I_1&=\frac{15}{64}\zeta(3)-\frac{\pi^2\log2}{32},\\
 I_2&=\frac{405}{1024}\zeta(5)-\frac{\pi^2\zeta(3)}{128}
                     -\frac{7\pi^4\log2}{1536},\\
 I_3&=\frac{7623}{16384}\zeta(7)-\frac{5\pi^2\zeta(5)}{2048}
       -\frac{7\pi^4\zeta(3)}{7680}-\frac{31\pi^6\log2}{61440}.
\end{align*}
\end{corollary}
\begin{proof}
Equation~\eqref{bip:eq:integral} with $(x,y)=(i,i)$ and $q=1$
gives $I_m=\Re Z_{2m,1}(i,i)$, since
$\Li_1(-t)=-\log(1+t)$ and $\Re(i/(1-it))=-t/(1+t^2)$.
In \eqref{bip:eq:R} use
$\mathcal B_{2k}(-1)=2\Li_{2k}(-1)$ and
$\mathcal B_{2k+1}(-1)=0$. Writing
$u_j=\Re\Li_j(i)$, the real part simplifies before any
zeta evaluation to
\[
 2I_m=-2\sum_{k=1}^{m}\Li_{2k}(-1)u_{w-2k}
           +w u_w-u_{2m}\log2-\Li_w(-1).
\]
Now $u_1=-\log2/2$ and
$u_j=-2^{-j}(1-2^{1-j})\zeta(j)$ for $j\ge2$.
Separate the $k=m$ term and substitute these expressions to obtain
\eqref{bip:eq:moment-family}.  Euler's even-zeta evaluations yield
the three displayed cases.
\end{proof}

Two weight-three integral identities also follow directly:
\begin{align*}
 \int_0^1\frac{x\log x\log(1+x)}{1+x^2}\,dx
  &=\frac{\pi^2\log2}{32}-\frac{15}{64}\zeta(3),\\
 \int_0^1\frac{x\Li_2(-x)}{1+x^2}\,dx
  &=\frac\pi4G-\frac{\pi^2\log2}{24}-\frac{33}{64}\zeta(3).
\end{align*}
These are ordinary real integrals with their constants fixed by the
convergent endpoint representation. They may be useful as simple
benchmark identities alongside the bicolored-polylogarithm form.

\subsection{Audit conclusion and what this does not settle}
The numeric Family--2 example can now be marked proved and replaced,
if desired, by \eqref{bip:eq:family2-short}; its older coordinates can
be preserved as an exact equivalent identity. No false equality was
found in that example. The correction is to its evidence status and
to an unnecessarily large choice of constants.

The master identity proves the real components in odd weight and the
imaginary components in even weight for every admissible Gaussian
mixed pair, and provides the same statement at all roots of unity.
It therefore supplies an exact replacement for searches in these
parity components. It does not prove nonreduction in the complementary
components.
In particular $S_6$ and $S_8$ have odd total weights seven and nine
and are imaginary components. They lie in the opposite parity sector;
this theorem proves neither frozen candidate. The rational
weight-nine vector already excluded in the source remains false, and
the newer $S_6,S_8$ candidates remain conjectural. The higher golden
ladders are also unaffected.

\paragraph{Attribution safeguard.}
When comparing formulas, Panzer uses increasing summation indices,
so both indices and arguments must be reversed to use his notation.
Panzer's footnote~16 explicitly flags a transposed product in
Nakamura's printed equation~(1.3): the product
$\zeta(a;-y)\zeta(b;x)$ must be replaced by
$\zeta(a;x)\zeta(b;-y)$ in that convention. The direct proof above
rebuilds the products from the three lattice regions and therefore
avoids relying on the misprinted equation. This is a literature
normalization issue, not a newly discovered source-manuscript error.

% Bibliography entries, if not already present in the host:
% \bibitem{bip:Panzer2017}
% E. Panzer, The parity theorem for multiple polylogarithms,
% Journal of Number Theory 172 (2017), 93--113.
% \url{https://arxiv.org/abs/1512.04482}.
% \bibitem{bip:Nakamura2012}
% T. Nakamura, A simple proof of the functional relation for the Lerch type
% Tornheim double zeta function, Tokyo Journal of Mathematics 35 (2012),
% 333--337. \url{https://arxiv.org/abs/1012.1144}.
